\chapter{Orientasi Mata Kuliah dan Dasar Arsitektur Enterprise}

\section*{Tujuan Pembelajaran}
\begin{itemize}
  \item Menjelaskan masalah organisasi yang dijawab EA (\textit{Enterprise Architecture}) dan nilai yang dapat diharapkan darinya, dengan merujuk bukti empiris.
  \item Mendeskripsikan sebuah organisasi memakai lima sudut pandang bersama: strategi, bisnis, data, aplikasi, dan teknologi.
  \item Menyusun piagam tim dan profil organisasi sebagai titik awal \textit{capstone}.
\end{itemize}

\section{Pendahuluan}

% TODO
% Cakupan: organisasi yang tumbuh lewat akuisisi dan proyek terpisah akhirnya punya sistem ganda,
%   data tersebar, dan keputusan TI tanpa gambaran utuh; pertanyaan pembuka: siapa yang memegang
%   peta besar organisasi ini?

\section{Enterprise, Arsitektur, dan Arsitektur Enterprise}

% TODO
% Cakupan: enterprise sebagai organisasi atau gabungan organisasi dengan tujuan bersama; arsitektur
%   sebagai struktur, hubungan, dan prinsip yang mengarahkan perubahan (ISO/IEC/IEEE 42010:2022 dan
%   TOGAF); EA sebagai praktik berkelanjutan, bukan sekadar dokumen; beda EA, arsitektur solusi, dan
%   arsitektur perangkat lunak.
% Kasus Medikara: lima rumah sakit hasil akuisisi dengan SIMRS berbeda versi sebagai gejala
%   organisasi yang tumbuh tanpa EA.

\section{Lima Sudut Pandang atas Organisasi}

% TODO
% Cakupan: strategi, bisnis, data, aplikasi, dan teknologi dengan contoh sehari-hari; bagaimana
%   perubahan di satu sudut pandang merambat ke sudut pandang lain; gambaran satu halaman sebagai
%   alat komunikasi dengan direksi.
% Kasus Medikara: satu kunjungan pasien rawat jalan dilihat dari kelima sudut pandang.

\section{Penyetaraan Istilah untuk Dua Kelompok Latar Belakang}

% TODO
% Cakupan: istilah teknis dasar (aplikasi, basis data, integrasi, API, cloud) dijelaskan dengan
%   analogi tanpa pemrograman; istilah bisnis dasar (strategi, model bisnis, model operasi, biaya,
%   nilai, pemangku kepentingan) untuk latar teknis; glosarium bersama tim sebagai alat kerja.
% Pendalaman teknis: hubungan dokumentasi arsitektur perangkat lunak dengan artefak EA menurut
%   Kotusev (2019).

\section{Nilai Arsitektur Enterprise dan Bukti Empirisnya}

% TODO
% Cakupan: mekanisme nilai EA menurut Tamm dkk. (2011); faktor keberhasilan pengelolaan EA menurut
%   Lange dkk. (2016); persepsi EA sebagai fungsi administratif TI menurut SAP LeanIX EA Insights
%   2025; penyebab umum kegagalan praktik EA.

\section{Peran Arsitek Enterprise}

% TODO
% Cakupan: tugas dan kompetensi arsitek enterprise; hubungan dengan direksi, manajer proyek, unit
%   bisnis, dan vendor; arsitek sebagai penerjemah antara bahasa bisnis dan bahasa teknologi.

\section{Proyek Capstone dan Kasus Medikara Health}

% TODO
% Cakupan: empat mata kuliah dalam satu proyek dan pembagian materinya; tonggak C1 sampai C4 beserta
%   jadwalnya; aturan izin, kerahasiaan, dan anonimisasi data bila memakai organisasi nyata; empat
%   peran tim yang bergilir (koordinator, analis bisnis, analis teknologi, pemodel).
% Kasus Medikara: profil, kondisi saat ini, strategi 2026 sampai 2030, dan pemicu perubahan;
%   ditegaskan sebagai organisasi fiktif.

\section{Aktivitas: Membentuk Tim dan Memotret Organisasi}

% TODO
% Tujuan aktivitas: membentuk tim campuran dan menghasilkan gambaran awal organisasi capstone.
% Peran dalam tim: koordinator menulis piagam tim; analis bisnis menyusun profil organisasi; analis
%   teknologi mendaftar sistem yang dipakai sehari-hari; pemodel menyusun glosarium bersama.
% Langkah: bentuk tim 4 sampai 5 orang berlatar campuran; pilih organisasi dan minta izin tertulis
%   atau pakai Medikara; tetapkan aturan anonimisasi; susun profil satu halaman; daftar sistem awal;
%   tulis piagam tim; sepakati 20 sampai 30 istilah dalam glosarium.
% Luaran dan tonggak capstone: piagam tim, profil organisasi, dan glosarium bersama (C1).
% Berkas sources yang direncanakan: sources/chapter01/templat-piagam-tim.md,
%   sources/chapter01/templat-profil-organisasi.md, sources/chapter01/panduan-izin-dan-anonimisasi.md,
%   sources/chapter01/glosarium-awal.md

\section{Ringkasan}

% TODO

\section{Latihan}

\begin{enumerate}
  \item Sebuah rumah sakit memakai tiga aplikasi pendaftaran pasien yang berbeda. Jelaskan dua kerugian bisnis dari kondisi itu dan cara EA membantu mencegahnya.
  \item Seorang manajer berkata, ``EA hanya menghasilkan diagram yang tidak dibaca orang.'' Tanggapi pernyataan itu dengan merujuk minimal satu temuan penelitian.
  \item Ambil satu contoh keputusan di sebuah organisasi yang gagal karena bagian bisnis dan bagian teknologi informasi memakai istilah dengan arti berbeda. Uraikan istilah yang disalahpahami dan dampaknya.
\end{enumerate}
